% Part: lambda-calculus
% Chapter: syntax
% Section: unique-readability

\documentclass[../../../include/open-logic-section]{subfiles}

\begin{document}

\olfileid[fr]{lam}{syn}{unq}
\olsection{Lecture unique}

Chaque terme possède-t-il une seule manière
d'être formé? Oui : chaque lambda-terme admet
une unique construction et, par conséquent,
une unique analyse syntaxique. Dans ce qui suit,
une \emph{formation} est un arbre décrivant la
construction d’un terme par une ou plusieurs
applications des règles de la \olref[trm]{defn:term}.

\begin{lem}\ollabel{lem:term-start}
Un terme commence soit par une variable,
soit par une parenthèse ouvrante.
\end{lem}

\begin{proof}
Une expression n'est un terme que si elle est
construite selon la \olref[trm]{defn:term}.
Si sa dernière étape est
\olref[trm]{defn:term-var}, elle est une variable.
Si sa dernière étape est
\olref[trm]{defn:term-abs} ou
\olref[trm]{defn:term-app}, elle commence
par une parenthèse ouvrante.
\end{proof}

\begin{lem}\ollabel{lem:app-start}
Le résultat d'une application commence soit
par deux parenthèses ouvrantes, soit par une
parenthèse ouvrante suivie d'une variable.
\end{lem}

\begin{proof}
Si $M$ résulte d'une application, il est
de la forme $(PQ)$ et commence donc par
une parenthèse ouvrante. Comme $P$ est un
terme, le \olref{lem:term-start} montre qu'il
commence à son tour soit par une parenthèse
ouvrante, soit par une variable.
\end{proof}

\begin{lem}\ollabel{lem:initial}
Aucun préfixe propre d'un terme n'est lui-même
un terme.
\end{lem}

\begin{prob}
Démontrer le \olref[lam][syn][unq]{lem:initial}
par induction sur la longueur des termes.
\end{prob}

\begin{prop}[Lecture unique] \ollabel{prop:unq}
Chaque terme possède une formation unique :
si une formation produit un terme~$M$,
aucune autre formation ne produit ce même terme.
\end{prop}

\begin{proof}
Nous raisonnons par induction sur la formation
des termes.
\begin{enumerate}
\item Supposons $M=x$, où $x$ est une variable.
  Les résultats d'une abstraction ou d'une
  application commencent toujours par une
  parenthèse; ces règles n'ont donc pas pu
  produire~$M$. Sa formation consiste en une
  seule application de la
  \olref[trm]{defn:term}\olref[trm]{defn:term-var}.
\item Supposons $M=(\lambd[x][N])$, où $x$
  est une variable et $N$ un terme. Ce n'est
  pas une variable isolée, donc la clause de la
  \olref[trm]{defn:term}\olref[trm]{defn:term-var}
  ne l'a pas produit. Ce n'est pas non plus
  le résultat d'une application, d'après le
  \olref{lem:app-start}. La dernière étape
  est donc nécessairement une abstraction
  sur~$N$. L'hypothèse d'induction garantit
  l'unicité de la formation de~$N$.
\item Supposons $M=(PQ)$, où $P$ et~$Q$
  sont des termes. Puisque $M$ commence par
  une parenthèse, il ne résulte pas de la
  \olref[trm]{defn:term}\olref[trm]{defn:term-var}.
  D'après le \olref{lem:term-start}, $P$ ne
  commence pas par~$\lambd$; $(PQ)$ n'est
  donc pas le résultat d'une abstraction.
  Supposons qu'une autre application construise
  le même terme sous la forme $(P'Q')$. Si
  $P\ne P'$, l'un des deux termes~$P$ et~$P'$
  est un préfixe propre de l'autre, ce qui
  contredit le \olref{lem:initial}. Ainsi $P=P'$
  et, par conséquent, $Q=Q'$. L'hypothèse
  d'induction assure alors l'unicité des
  formations de~$P$ et de~$Q$.
\end{enumerate}
\end{proof}

La proposition peut se reformuler plus
simplement comme suit.
\begin{prop}
Un terme~$M$ possède exactement l'une des
formes suivantes :
\begin{enumerate}
\item $x$, où la variable~$x$ est déterminée
  de manière unique par~$M$;
\item $(\lambd[x][N])$, où la variable~$x$
  et le terme~$N$ sont tous deux déterminés
  de manière unique par~$M$;
\item $(PQ)$, où les termes~$P$ et~$Q$ sont
  déterminés de manière unique par~$M$.
\end{enumerate}
\end{prop}

\end{document}
